% Phase 6 only: qualified-capacity supply, demand, and scalar market clearing.
\section{Qualified Manufacturing-Service Market}
\label{sec:p06-cmo-market}

This section closes only the market for qualified manufacturing-service
capacity. Developers and qualified suppliers take a conjectured price \(p_m\)
as given when making their individual choices. The market-consistent price
\(p_m^*\) is then selected by one scalar clearing equation. No entry, labor,
capital, product-market, or welfare market is added.

\subsection{Qualified-capacity supply}
\label{subsec:p06-supply}

Qualified supplier \(j\), with predetermined efficiency \(z_j\), chooses
capacity \(s_j\geq0\) to solve
\begin{equation}
  \max_{s_j\geq0}
  \left\{
    p_ms_j-\Psi(s_j;z_j)
  \right\}.
  \label{eq:p06-supplier-problem}
\end{equation}
Capacity has units \(\mathsf B\) per decision cohort, \(p_m\) has units
\(\mathsf C/\mathsf B\), and \(\Psi\) has units \(\mathsf C\).

For every \(z_j\), assume
\[
  \Psi(0;z_j)=0,\qquad \Psi_s(0;z_j)=0,
\]
\[
  \Psi_s>0\ \text{for }s_j>0,\qquad
  \Psi_{ss}>0,\qquad
  \Psi_{sz}<0,
\]
and \(\Psi_s(s_j;z_j)\to+\infty\) as \(s_j\to+\infty\).
At \(p_m=0\), the unique optimum is \(s_j^*=0\). For \(p_m>0\), the
solution is interior and satisfies
\begin{equation}
  p_m=\Psi_s\!\left(s_j^*(p_m,z_j);z_j\right).
  \label{eq:p06-supplier-capacity}
\end{equation}
The objective's second derivative is
\(-\Psi_{ss}<0\), so the capacity choice is unique. The
implicit-function theorem gives
\begin{equation}
  \frac{\partial s_j^*}{\partial p_m}
  =
  \frac{1}{\Psi_{ss}}>0,
  \qquad
  \frac{\partial s_j^*}{\partial z_j}
  =
  -\frac{\Psi_{sz}}{\Psi_{ss}}>0.
  \label{eq:p06-capacity-derivatives}
\end{equation}

Aggregate qualified-capacity supply is
\begin{equation}
  S_m(p_m)
  =
  \int s_j^*(p_m,z)\,dH_C(z).
  \label{eq:p06-aggregate-supply}
\end{equation}
Under the stated integrability conditions, \(S_m\) is continuous and
strictly increasing, \(S_m(0)=0\), and
\(S_m(p_m)\to+\infty\) as \(p_m\to+\infty\). Supply technology and
\(H_C\) are invariant to \(M\); MAH creates no direct supply shift.

\subsection{Entrusted-capacity demand per planning-stage project}
\label{subsec:p06-project-demand}

At a candidate price \(p_m\), define developer \(i\)'s expected qualified
capacity per planning-stage project as
\begin{equation}
  \chi_i^E(p_m;M)
  =
  \int
  b(m)
  \mathbf 1
  \left\{
    r_i^*(q,m;M,p_m)=E
  \right\}
  dF(q,m).
  \label{eq:p06-entrusted-capacity}
\end{equation}
This is an integral of deterministic route choices, not a probabilistic
route-share formula.

The entrusted value falls with price at slope \(-b(m)<0\), while the other
route values do not contain \(p_m\). Consequently, for each project the set
of prices at which \(E\) is chosen is an interval from below: once a higher
price makes \(E\) lose to another route, a still higher price cannot restore
it. Hence \(\chi_i^E(p_m;M)\) is weakly decreasing in \(p_m\).

Away from the measure-zero set of route ties, the envelope derivative of the
optimized project value is
\[
  \frac{\partial W_i(q,m;M,p_m)}{\partial p_m}
  =
  -b(m)
  \mathbf 1\{r_i^*(q,m;M,p_m)=E\}.
\]
Dominated convergence then gives
\begin{equation}
  \frac{\partial\Omega_i(M,p_m)}{\partial p_m}
  =
  -\chi_i^E(p_m;M)
  \leq0
  \quad\text{almost everywhere in }p_m.
  \label{eq:p06-omega-price-envelope}
\end{equation}
For \(\Omega_i>0\), the Phase~5 optimizer therefore satisfies
\begin{equation}
  \frac{\partial x_i^*(M,p_m)}{\partial p_m}
  =
  -\frac{x_i^*(M,p_m)}
  {\nu\Omega_i(M,p_m)}
  \chi_i^E(p_m;M)
  \leq0.
  \label{eq:p06-advancement-price-response}
\end{equation}
At the zero-value corner, \(x_i^*=0\); globally the optimal advancement
intensity remains weakly decreasing in \(p_m\).

\subsection{Aggregate demand}
\label{subsec:p06-demand}

Developer \(i\)'s expected qualified-capacity demand is
\[
  a_ix_i^*(M,p_m)\chi_i^E(p_m;M).
\]
Aggregating over developer heterogeneity gives study-related demand
\begin{equation}
  D_m^{\mathrm{MAH}}(p_m;M)
  =
  \int
  a_ix_i^*(M,p_m)\chi_i^E(p_m;M)
  dH(a,k).
  \label{eq:p06-study-demand}
\end{equation}
This expression contains both price feedbacks required by the model:
\(x_i^*\) responds through expected project value, and
\(\chi_i^E\) responds through deterministic route selection.

Both factors are nonnegative and weakly decreasing in \(p_m\), so
\(D_m^{\mathrm{MAH}}\) is weakly decreasing. Individual route indicators can
jump at a cutoff. Aggregate demand is nevertheless continuous: when
\(p_m^n\to p_m\), route indicators converge for every project--developer
pair except the zero-measure tie set, \(x_i^*\) is continuous through
\(\Omega_i\), and the common integrable envelope permits dominated
convergence. Thus aggregate regularity comes from continuous heterogeneity,
not smooth random utility.

Let \(D_m^B(p_m)\) be finite, continuous, nonnegative, weakly decreasing
background demand with
\[
  D_m^B(0)>0,
  \qquad
  \lim_{p_m\to\infty}D_m^B(p_m)=0.
\]
It is policy invariant and allows a CMO market to exist before MAH. Total
demand is
\begin{equation}
  D_m(p_m;M)
  =
  D_m^B(p_m)
  +
  D_m^{\mathrm{MAH}}(p_m;M).
  \label{eq:p06-total-demand}
\end{equation}
It is finite, continuous, nonnegative, and weakly decreasing.

\subsection{Market clearing, existence, and uniqueness}
\label{subsec:p06-clearing}

The qualified-capacity market clears at
\begin{equation}
  D_m(p_m^*;M)=S_m(p_m^*).
  \label{eq:p06-market-clearing}
\end{equation}
At zero price,
\begin{equation}
  D_m(0;M)\geq D_m^B(0)>0=S_m(0).
  \label{eq:p06-existence-boundaries}
\end{equation}
As \(p_m\to\infty\), \(W_i^E\to-\infty\) pointwise because
\(b(m)>0\), so \(\chi_i^E\to0\) and study-related demand tends to zero
under the integrable envelope. Background demand also tends to zero, whereas
aggregate supply tends to infinity. Continuity therefore gives at least one
positive clearing price.

Total demand is weakly decreasing and aggregate supply is strictly
increasing. Their difference is strictly decreasing, so the clearing price
is unique.

\subsection{Closed scalar solution order}
\label{subsec:p06-solution-order}

For each regime \(M\), the equilibrium can be computed without an unresolved
cycle:
\begin{enumerate}
  \item conjecture \(p_m\);
  \item evaluate route values and deterministic \(r_i^*(q,m;M,p_m)\);
  \item integrate \(\Omega_i(M,p_m)\) and compute \(x_i^*(M,p_m)\);
  \item integrate \(\chi_i^E\) and \(D_m^{\mathrm{MAH}}\);
  \item solve supplier capacity and aggregate \(S_m(p_m)\);
  \item find the unique zero of \(D_m(p_m;M)-S_m(p_m)\).
\end{enumerate}
Every object on the right side is a function of the candidate price, leaving
one scalar equation in \(p_m\).

When \(M=0\), route \(E\) is unavailable, so
\(D_m^{\mathrm{MAH}}(p_m;0)=0\); background demand still supports the market.
When \(M=1\), study demand is nonnegative and supply is unchanged. Hence the
post-MAH clearing price is weakly above the pre-MAH price, and it is strictly
higher if study demand is positive at the pre-MAH price. This is an endogenous
scarcity-price response, not a direct policy shift in \(p_m^*\) or CMO
technology.

\subsection{Scope audit}
\label{subsec:p06-scope}

Phase 6 closes qualified manufacturing-service capacity only. It introduces
no supplier entry, no direct policy supply shift, no logit route smoothing,
and no additional market-clearing condition. The price response is generated
jointly by supplier optimization, deterministic route selection, project
advancement, and scalar market clearing.
